\section{Introduction}
\label{sec:intro}

Retrieval-augmented generation~\cite{lewis2020rag} allows a language model to answer questions
from a document it was not trained on. For a single PDF, the usual approach is simple. The text is
split into chunks of a fixed number of characters, each chunk is embedded with a sentence encoder
such as Sentence-BERT~\cite{reimers2019sbert}, the chunks nearest to the question are found in a
vector index, and these are sent to the model together with the question. This approach is used
in many tutorials, course projects and portfolio repositories. A common configuration, which is
also the one in our app, is 1,000-character chunks, the MiniLM encoder~\cite{wang2020minilm}, a
FAISS index~\cite{johnson2021faiss} and the top four chunks. Such apps are rarely tested on a
public benchmark. When they are, the report usually covers either retrieval or the answers, and
does not show where between the two the answers are lost.

Research systems for scientific question answering are evaluated more carefully, but they are
built differently. PaperQA~\cite{lala2023paperqa} and PaperQA2~\cite{skarlinski2024paperqa2}
search the literature, summarise each retrieved passage for the question before answering, and
may decline to answer. OpenScholar~\cite{asai2024openscholar} retrieves from tens of millions of
papers and revises its answers using feedback. These systems answer questions across the
literature, whereas the small apps answer questions about a paper the user already has. This is
the setting of the QASPER benchmark~\cite{dasigi2021qasper}, in which questions were written by
readers who had seen only the title and abstract of an NLP paper, and answered by other
annotators from the full text, who also marked the evidence paragraphs.

We built one such app, CitePage, which answers questions about uploaded PDFs through the Groq
API, and in this paper we run its own pipeline over the QASPER development set. We first wanted
to know how the app behaved as we had released it, and how often retrieval, rather than the
model, was what kept an answer from being right. We then tried the settings people usually
change (chunk size, embedding model, keyword search and reranking) while keeping the prompt the
same length. Users also load several papers at once, so we checked whether retrieval still finds
the paper a question is about. Finally, we checked the measurements themselves. We graded a
sample of answers by hand to see how far to trust the language-model judge, asked every question
three times to see how much the results move, and looked at whether the pages an answer cites
contain what it says.

These results come from one app, one benchmark of NLP papers and one answering model, and we
do not propose a new method. Instead, we show where a typical pipeline loses its answers, and
which of the changes we tried helped and which did not. Every number in the paper can be traced
to the per-question outputs we release.
